\documentclass[10pt,a4paper]{amsart}
\usepackage[margin=19mm]{geometry}
\usepackage{amsmath,amssymb,mathtools}
\usepackage[scr=rsfs,cal=euler]{mathalfa}
\usepackage{xfrac}
\usepackage[unicode,hidelinks,bookmarks=false]{hyperref}
\newcommand{\ie}{i.e.\ }
\newcommand{\cf}{cf.\ }
\newcommand{\resp}{respectively\ }
\newcommand{\Subsection}{\subsection}
\providecommand{\nobreakdash}{\nobreak\hbox{-}}
\setlength{\emergencystretch}{2em}
\multlinegap0pt
\usepackage{amssymb}
\usepackage{xcolor}
\newtheorem{lemma}{Lemma}
\title{Additive problems on $\lfloor p^c \rfloor$}
\author{Lingyu Guo, Victor Zhenyu Guo, Li Lu}
\date{Source: arXiv:2505.19833v2 (CC BY 4.0)}
\begin{document}
\maketitle

\noindent\emph{Source and licence.} Additive problems on $\lfloor p^c \rfloor$. Version arXiv:2505.19833v2 (CC BY 4.0), distributed by arXiv under the Creative Commons Attribution 4.0 International licence, http://creativecommons.org/licenses/by/4.0/, as recorded in the arXiv licence metadata for this exact version and shown on the abstract page https://arxiv.org/abs/2505.19833v2. Full licence notice: \texttt{licenses/arxiv-2505.19833-CC-BY-4.0.txt}.\par\smallskip

\noindent\textit{Editorial scope.} Counted unit: the article's lemma lem:mean\_value (source0.tex lines 470--484). The article's complete original proof of it is delivered with it (source0.tex lines 486--488, one \texttt{proof} environment of 3 lines). No mathematics is written, repaired or re-derived: the statement and the proof are the article's own lines, in the article's own notation, and no retained line is substituted or re-flowed.  No further source-local context is needed: every cross-reference of every retained line resolves inside the two delivered spans.  Nothing was omitted from any retained span, so the package carries no omission record.  The retained lines cite 1 works, whose entries are transcribed verbatim from the article's own bibliography source in the e-print and delivered with short editorial labels recorded in the item's reference mapping.  No retained line contains a figure environment, an image-inclusion command or drawing code, so no image is delivered and the package declares no asset dependency.  Editorial formatting only, in addition: an AMS \texttt{amsart} preamble that restates the article's own declarations and macros used by the retained lines, namely the environments \texttt{lemma} on separate counters, together with the standard packages those lines need.  The retained environments print in this document as Lemma 1 in order of appearance, whereas the article prints the same items with the numbering of its own full document.  The complete native source of the article as distributed in the arXiv e-print for this exact version is bundled unchanged as \texttt{sources/178774/source0.tex}.
\begin{lemma}
\label{lem:mean_value}
Let 
$$
S(s,\chi)=\sum_{n=1}^{N}a_n\chi(n)n^{-s},
$$
and let $\mathcal{A}_\chi$ be a finite set of complex numbers $s=\sigma+it$. Let $T_0,T,\sigma_0,\delta$ be real numbers such that for each $\chi$,
$$
T_0+\frac{\delta}{2}\leqslant t\leqslant T_0+T-\frac{\delta}{2},\quad \sigma\geqslant \sigma_0 \quad \text{and} \quad |t-t'|\geqslant \delta 
$$
for all $s,s'\in\mathcal{A}_\chi$ and $s\neq s'$. Then 
$$
\sum_{\chi}\sum_{s\in \mathcal{A}_\chi}|S(s,\chi)|^2\ll (qT+N)(\delta^{-1}+\log N)\sum_{n=1}^{N}|a_n|^2n^{-2\sigma_0}\bigg( 1+\log \frac{\log 2N}{\log 2n} \bigg).
$$ 
\end{lemma}

\begin{proof}
See \cite[Theorem~7.6]{Montgomery2}.
\end{proof}

\begin{thebibliography}{99}
\bibitem[Montgo]{Montgomery2} H.~L.~Montgomery, Topics in multiplicative number theory. Lecture Notes in Mathematics, Vol. 227. Springer--Verlag, Berlin-New York, 1971.
\end{thebibliography}
\end{document}
